% ============================================================
\chapter{Expectation, Variance, Moments}
\label{ch:expectation}
% ============================================================

If you play a game of chance a large number of times, your average gain per game will converge to a precise number: the \emph{expectation}. This result, anticipated by Pascal and Fermat as early as 1654, is the law of large numbers. Expectation summarises the ``location'' of a random variable; \emph{variance}, introduced by Ronald Fisher in 1918, measures its ``spread.'' Together, these two parameters capture the essence of a distribution's behaviour, and higher-order moments---skewness, kurtosis---refine the portrait.

\begin{intuition}
Expectation and variance are the two fundamental summary parameters of a random
variable: location and spread. This chapter defines them rigorously, establishes
their properties, and introduces higher-order moments.
\end{intuition}

% ------------------------------------------------------------
\section{Expectation}
% ------------------------------------------------------------

\begin{definition}[Expectation --- discrete case]
Let $X$ be a discrete random variable taking values $x_1, x_2, \ldots$
The \textbf{expectation} (or \textbf{mean}) of $X$ is:
\[
    \E[X] = \sum_{k} x_k\,\PP(X = x_k),
\]
provided $\sum_k \abs{x_k}\,\PP(X = x_k) < \infty$ (absolute convergence).
\end{definition}

\begin{definition}[Expectation --- continuous case]
If $X$ has density $f_X$, the \textbf{expectation} is:
\[
    \E[X] = \int_{-\infty}^{+\infty} x\,f_X(x)\,dx,
\]
provided $\int \abs{x}\,f_X(x)\,dx < \infty$.
\end{definition}

\begin{theorem}[Law of the unconscious statistician (LOTUS)]
\label{thm:lotus}
If $g : \R \to \R$ is measurable:
\begin{itemize}
    \item \textbf{Discrete:}
    $\E[g(X)] = \sum_k g(x_k)\,\PP(X = x_k)$.
    \item \textbf{Continuous:}
    $\E[g(X)] = \int_{-\infty}^{+\infty} g(x)\,f_X(x)\,dx$.
\end{itemize}
(Provided the sums/integrals converge absolutely.)
\end{theorem}

\begin{intuition}[title=\textbf{Why LOTUS is so useful}]
The transfer formula lets us compute $\E[g(X)]$ directly from the distribution of
$X$, \emph{without} having to find the distribution of $Y = g(X)$.
\end{intuition}

\begin{theorem}[Properties of expectation]
\label{thm:exp_properties}
\begin{enumerate}[label=(\roman*)]
    \item \textbf{Linearity:} $\E[aX + bY] = a\E[X] + b\E[Y]$ for all
          $a, b \in \R$.
    \item \textbf{Positivity:} if $X \geq 0$ a.s., then $\E[X] \geq 0$.
    \item \textbf{Monotonicity:} if $X \leq Y$ a.s., then $\E[X] \leq \E[Y]$.
    \item $\abs{\E[X]} \leq \E[\abs{X}]$.
    \item \textbf{Product for independent variables:}
          if $X \indep Y$ and $\E[\abs{X}], \E[\abs{Y}] < \infty$, then
          $\E[XY] = \E[X]\,\E[Y]$.
\end{enumerate}
\end{theorem}

\begin{example}
Let $X \sim \mathrm{Exp}(\lambda)$. Then:
\[
    \E[X] = \int_0^{\infty} x\,\lambda e^{-\lambda x}\,dx
    = \left[-x\,e^{-\lambda x}\right]_0^{\infty}
    + \int_0^{\infty} e^{-\lambda x}\,dx
    = 0 + \frac{1}{\lambda} = \frac{1}{\lambda}.
\]
\end{example}

% ------------------------------------------------------------
\section{Variance}
% ------------------------------------------------------------

\begin{definition}[Variance and standard deviation]
The \textbf{variance} of $X$ (when $\E[X^2] < \infty$) is:
\[
    \Var(X) = \E\bigl[(X - \E[X])^2\bigr].
\]
The \textbf{standard deviation} is $\sigma(X) = \sqrt{\Var(X)}$.
\end{definition}

\begin{theorem}[K\"onig--Huygens formula]
\label{thm:konig}
\[
    \Var(X) = \E[X^2] - (\E[X])^2.
\]
\end{theorem}

\begin{proof}
\begin{align*}
    \Var(X) &= \E[(X - \E[X])^2]
    = \E[X^2 - 2X\E[X] + (\E[X])^2] \\
    &= \E[X^2] - 2\E[X]\E[X] + (\E[X])^2
    = \E[X^2] - (\E[X])^2. \qedhere
\end{align*}
\end{proof}

\begin{proposition}[Properties of variance]
\begin{enumerate}[label=(\roman*)]
    \item $\Var(X) \geq 0$, with equality iff $X$ is constant a.s.
    \item $\Var(aX + b) = a^2\Var(X)$.
    \item If $X \indep Y$, then $\Var(X + Y) = \Var(X) + \Var(Y)$.
    \item In general: $\Var(X+Y) = \Var(X) + \Var(Y) + 2\Cov(X,Y)$.
\end{enumerate}
\end{proposition}

\begin{keyformulas}[title=\textbf{Summary: expectation and variance}]
\[
    \E[aX+b] = a\E[X]+b, \qquad \Var(aX+b) = a^2\Var(X), \qquad
    \Var(X) = \E[X^2] - \E[X]^2.
\]
\end{keyformulas}

% ------------------------------------------------------------
\section{Higher-Order Moments}
% ------------------------------------------------------------

\begin{definition}[Moment of order $r$]
The \textbf{$r$-th moment} of $X$ ($r \in \N^*$) is:
\[
    m_r = \E[X^r],
\]
when this expectation exists. The \textbf{$r$-th central moment} is
$\mu_r = \E[(X - \E[X])^r]$.
\end{definition}

\begin{definition}[Skewness and kurtosis]
\begin{itemize}
    \item The \textbf{skewness} is $\gamma_1 = \mu_3/\sigma^3$.
    \item The \textbf{excess kurtosis} is $\gamma_2 = \mu_4/\sigma^4 - 3$.
\end{itemize}
For the normal distribution, $\gamma_1 = 0$ and $\gamma_2 = 0$.
\end{definition}

\begin{remark}
$\gamma_1 > 0$ indicates a heavier right tail; $\gamma_1 < 0$ a heavier left tail.
$\gamma_2 > 0$ indicates tails heavier than the normal (leptokurtic distribution).
\end{remark}

% ------------------------------------------------------------
\section{Fundamental Inequalities}
% ------------------------------------------------------------

\begin{theorem}[Markov's inequality]
\label{thm:markov_ineq}
If $X \geq 0$ and $a > 0$:
\[
    \PP(X \geq a) \leq \frac{\E[X]}{a}.
\]
\end{theorem}

\begin{proof}
$\E[X] = \E[X\,\mathbf{1}_{X \geq a}] + \E[X\,\mathbf{1}_{X < a}]
\geq \E[X\,\mathbf{1}_{X \geq a}] \geq a\,\PP(X \geq a)$.
\end{proof}

\begin{theorem}[Chebyshev's inequality]
\label{thm:chebyshev}
If $\Var(X)$ exists:
\[
    \PP\bigl(\abs{X - \E[X]} \geq \varepsilon\bigr)
    \leq \frac{\Var(X)}{\varepsilon^2}, \quad \forall\,\varepsilon > 0.
\]
\end{theorem}

\begin{proof}
Apply Markov's inequality to $Y = (X - \E[X])^2$ with $a = \varepsilon^2$:
\[
    \PP(\abs{X - \E[X]} \geq \varepsilon) = \PP(Y \geq \varepsilon^2)
    \leq \frac{\E[Y]}{\varepsilon^2} = \frac{\Var(X)}{\varepsilon^2}. \qedhere
\]
\end{proof}

\begin{theorem}[Jensen's inequality]
\label{thm:jensen}
If $\varphi$ is convex and $\E[\abs{X}] < \infty$:
\[
    \varphi(\E[X]) \leq \E[\varphi(X)].
\]
If $\varphi$ is concave, the inequality is reversed.
\end{theorem}

\begin{example}
Taking $\varphi(x) = x^2$ (convex), we recover $(\E[X])^2 \leq \E[X^2]$,
i.e.\ $\Var(X) \geq 0$.
\end{example}

\begin{theorem}[Cauchy--Schwarz inequality]
\label{thm:cauchy_schwarz}
\[
    \abs{\E[XY]}^2 \leq \E[X^2]\,\E[Y^2].
\]
Equality holds if and only if $Y = aX + b$ a.s.\ for some $a, b \in \R$.
\end{theorem}

% ------------------------------------------------------------
\section{Conditional Expectation (Introduction)}
% ------------------------------------------------------------

\begin{definition}[Discrete conditional expectation]
If $X$ and $Y$ are discrete:
\[
    \E[X \mid Y = y] = \sum_x x\,\PP(X = x \mid Y = y).
\]
\end{definition}

\begin{theorem}[Law of total expectation]
\[
    \E[X] = \E[\E[X \mid Y]] = \sum_y \E[X \mid Y = y]\,\PP(Y = y).
\]
\end{theorem}

\begin{theorem}[Law of total variance]
\[
    \Var(X) = \E[\Var(X \mid Y)] + \Var(\E[X \mid Y]).
\]
\end{theorem}

\begin{example}
Let $N \sim \mathrm{Poisson}(\lambda)$ and, conditionally on $N = n$,
$S = X_1 + \cdots + X_n$ where the $X_i$ are i.i.d.\ with $\E[X_i] = \mu$
and $\Var(X_i) = \sigma^2$. Then:
\begin{align*}
    \E[S] &= \E[\E[S \mid N]] = \E[N\mu] = \lambda\mu, \\
    \Var(S) &= \E[\Var(S \mid N)] + \Var(\E[S \mid N])
    = \E[N\sigma^2] + \Var(N\mu) = \lambda\sigma^2 + \lambda\mu^2.
\end{align*}
\end{example}

% ------------------------------------------------------------
\section{Exercises}
% ------------------------------------------------------------

\begin{exercise}
Compute $\E[X]$ and $\Var(X)$ for $X \sim \mathcal{N}(\mu, \sigma^2)$ directly
by integration.
\end{exercise}

\begin{exercise}
Let $X$ have density $f(x) = 6x(1-x)\,\mathbf{1}_{[0,1]}(x)$. Compute $\E[X]$,
$\E[X^2]$, $\Var(X)$, $\gamma_1$, and $\gamma_2$.
\end{exercise}

\begin{exercise}
Prove Markov's inequality. Deduce that if $\E[e^{tX}]$ exists for $t > 0$, then
$\PP(X \geq a) \leq e^{-ta}\,\E[e^{tX}]$ for all $a$ (Chernoff bound).
\end{exercise}

\begin{exercise}
Let $X$ be uniformly distributed on $\{1, 2, \ldots, n\}$. Compute $\E[X]$,
$\E[X^2]$, and $\Var(X)$.
\end{exercise}

\begin{exercise}
Prove the Cauchy--Schwarz inequality by considering the polynomial
$p(t) = \E[(X + tY)^2] \geq 0$ and studying its discriminant.
\end{exercise}

\begin{exercise}
A fair die is rolled $N$ times, where $N \sim \mathrm{Poisson}(10)$.
Let $S$ be the sum of the results. Compute $\E[S]$ and $\Var(S)$.
\end{exercise}

\begin{exercise}
Let $X$ be a non-negative random variable. Show that
\[
    \E[X] = \int_0^{\infty} \PP(X > t)\,dt.
\]
\textit{Hint:} interchange the integral and the expectation.
\end{exercise}

\begin{exercise}
Show that $\Var(X) = \min_{c \in \R}\E[(X-c)^2]$, with the minimum attained
at $c = \E[X]$.
\end{exercise}
